\documentclass[12pt]{article}
\usepackage[utf8]{inputenc}
\usepackage{amsmath, amssymb, amsthm}
\usepackage{geometry}
\usepackage{booktabs}
\geometry{margin=1in}

\title{Problem 45: Triangular, Pentagonal, and Hexagonal}
\author{Project Euler Solutions}
\date{}

\newtheorem{theorem}{Theorem}

\begin{document}
\maketitle

\section{Problem Statement}

Triangle, pentagonal, and hexagonal numbers are given by:
\begin{align*}
T_n &= \frac{n(n+1)}{2}, \\
P_n &= \frac{n(3n-1)}{2}, \\
H_n &= n(2n-1).
\end{align*}

We know that $T_{285} = P_{165} = H_{143} = 40755$. Find the next number that is simultaneously triangular, pentagonal, and hexagonal.

\section{Key Theorem}

\begin{theorem}
Every hexagonal number is a triangle number. Specifically, $H_m = T_{2m-1}$ for all $m \geq 1$.
\end{theorem}

\begin{proof}
\[
T_{2m-1} = \frac{(2m-1)(2m-1+1)}{2} = \frac{(2m-1)(2m)}{2} = m(2m-1) = H_m. \qedhere
\]
\end{proof}

\section{Reduced Problem}

Since every hexagonal number is automatically triangular, we only need to find the next number that is both \textbf{hexagonal} and \textbf{pentagonal}. That is, we seek the smallest $m > 143$ such that $H_m$ is also pentagonal.

\section{Pentagonal Test}

A positive integer $v$ is pentagonal if and only if
\[
n = \frac{1 + \sqrt{1 + 24v}}{6}
\]
is a positive integer. Equivalently, $1 + 24v$ must be a perfect square $s^2$ with $s \equiv 1 \pmod{6}$.

\section{Connection to Pell's Equation}

Setting $H_m = P_k$:
\[
m(2m-1) = \frac{k(3k-1)}{2} \implies 2m(2m-1) = k(3k-1).
\]

Let $x = 4m - 1$ and $y = 6k - 1$. Then:
\[
\frac{(x+1)}{2} \cdot \frac{x}{1} = \frac{(y+1)}{6} \cdot \frac{(y+2)}{2}
\]

After algebraic manipulation, this reduces to a generalized Pell equation:
\[
2x^2 - 3y^2 = -1
\]
where $x = 4m-1$ and $y = 6k-1$. The fundamental solution gives rise to a recurrence generating all solutions. The solution after $m = 143$ is $m = 27693$.

\section{Answer}

\[
\boxed{1533776805}
\]
\section{Complexity}

\begin{itemize}
    \item \textbf{Time:} $O(M)$ where $M = 27693$ is the hexagonal index. Each pentagonal test is $O(1)$ using the square root formula.
    \item \textbf{Space:} $O(1)$.
\end{itemize}

\end{document}
